\documentclass{article}
\usepackage[OT4]{fontenc}
\usepackage[polish]{babel}
\usepackage[legalpaper, margin=2cm]{geometry}
\usepackage{mathtools}
\usepackage{float}
\usepackage{hyperref}
\usepackage{placeins}
\usepackage{chngcntr}
\usepackage{graphicx}


\newif\ifszybkakompilacja
\szybkakompilacjatrue
% \szybkakompilacjatrue

\newcommand{\raportimg}[1]{%
    \centering
    \ifszybkakompilacja
        \fbox{%
            \parbox[c][0.24\textheight][c]{0.78\textwidth}{%
                \centering
                \texttt{#1}\\[0.4em]
                \small miejsce na wykres
            }%
        }%
    \else
        \includegraphics[width=0.78\textwidth,height=0.32\textheight,keepaspectratio]{#1}%
    \fi
}


% ============================================================
% TRYB WYKRESÓW Z SEKCJI 6
% ============================================================
% Jeśli Overleaf Free łapie timeout, zostaw:
% \pokazsekcjaszescfalse
%
% Wtedy sekcja 6 pokazuje ramki z nazwami plików.
% Po zmniejszeniu PNG do 100--120 DPI możesz zmienić na:
% \pokazsekcjaszesctrue
% ============================================================

\newif\ifpokazsekcjaszesc
\pokazsekcjaszesctrue
% \pokazsekcjaszesctrue

\newcommand{\wykrespngwniosek}[4]{%
\begin{figure}[H]
    \centering
    \ifpokazsekcjaszesc
        \includegraphics[width=0.72\textwidth,height=0.24\textheight,keepaspectratio]{#1}%
    \else
        \fbox{%
            \parbox[c][0.20\textheight][c]{0.72\textwidth}{%
                \centering
                \texttt{#1}\\[0.4em]
                \small miejsce na wykres
            }%
        }%
    \fi
    \caption{#2}
    \label{fig:#3}
\end{figure}
\noindent\textbf{Wniosek:} #4
\par\vspace{0.2cm}
}

\title{Sprawozdanie z pierwszej połowy projektu\\
Przedmiotu "Wizualizacja Danych"}
\author{Zespół 4 - Jan Domański, Maciej Kurzydłowski, Jan Machowski}
\date{sem. 26L}

\begin{document}
\maketitle
\tableofcontents
\newpage

% ============================================================
\section{Standarized price-to-income ratio}
% ============================================================

\subsection{Źródło danych}
\href{https://ec.europa.eu/eurostat/databrowser/view/tipsho60/default/table?lang=en}{Eurostat tipsho60 - Standarized price-to-income ratio - annual data}\\
\url{https://ec.europa.eu/eurostat/databrowser/view/tipsho60/default/table?lang=en}

\subsection{Opis zbioru i istotnych pojęć}
\begin{itemize}
    \item Zbiór przedstawia uśrednione wartości PTI dla poszczególnych krajów;
    \item Dane zostały zebrane na przestrzeni lat 2000--2024;
    \item Dane pochodzą z 27 krajów członkowskich Unii Europejskiej;
    \item \textbf{AGDI} - Average Gross Domestic Income - skorygowany dochód do dyspozycji brutto gospodarstwa domowego;
    \item \textbf{PTI} - Price-to-income, ustandaryzowany stosunek średniej ceny mieszkania do AGDI gospodarstwa domowego;
    \item \textbf{RoC} - Rate of Change, współczynnik zmiany wartości PTI względem poprzedniego roku;
    \item Większość danych z okresu 2000--2004 jest estymatami, dlatego należy interpretować je ostrożnie;
    \item Zwiększenie lub zmniejszenie indeksu PTI nie oznacza jednoznacznie zmiany ceny mieszkania, lecz zmianę relacji ceny mieszkania do dochodu;
    \item Indeks PTI można interpretować jako średnią względną wartość domu w danym kraju.
\end{itemize}

\subsection{Wizualizacje zbioru}

\subsubsection{Statystyka opisowa Price-to-Income}
\begin{figure}[H]
    \raportimg{phts/1/png/1.1.png}
    \caption{Statystyka opisowa danych PTI względem kraju.}
\end{figure}

Wykres przedstawia podstawowe statystyki opisowe wartości PTI względem średniej długoterminowej dla zbioru krajów.
Średnia na przestrzeni lat prezentuje się stosunkowo równomiernie, a większość krajów znajduje się w zakresie około 90\%--120\% PTI.
Widać jednak, że odchylenie jednego z państw -- Rumunii -- dominuje nad pozostałymi.
Sugeruje to, że dane dla tego kraju są szczególnie zmienne, co również widać po wartościach minimalnych i maksymalnych.

\subsubsection{Indeks Price-to-Income względem kraju}
\begin{figure}[H]
    \raportimg{phts/1/png/1.2.png}
    \caption{Wartość indeksu PTI względem kraju na przestrzeni lat względem roku 2015.}
\end{figure}

Dane przedstawione na wykresie są wartościami PTI względem roku 2015 dla poszczególnych krajów, stąd kumulacja linii w tym punkcie.
Dla klarowności dane podzielono na Europę Wschodnią i Europę Zachodnią.
Najbardziej zmienna jest Rumunia, uzyskująca szczyt w 2007 roku powyżej 250\% PTI, a następnie stopniowo spadająca.
Na wykresie widać także tendencję spadkową indeksów PTI wielu krajów w latach 2007--2015, co można wiązać z kryzysem na rynku mieszkaniowym po 2008 roku.

\subsubsection{Price-to-Income w 2024 względem średniej długoterminowej}
\begin{figure}[H]
    \raportimg{phts/1/png/1.4.png}
    \caption{Indeks PTI krajów w 2024 roku względem średniej długoterminowej.}
\end{figure}

Wykres przedstawia różnicę najnowszej wartości indeksu PTI krajów europejskich względem średniej długoterminowej.
Najwyższe dodatnie odchylenia widoczne są między innymi dla Portugalii, Luksemburga i Austrii.
Warto zauważyć, że podział na Europę Wschodnią oraz Zachodnią nie determinuje jednoznacznie tego, czy mieszkania są relatywnie bardziej wartościowe względem dochodu.

\subsubsection{Zmiana Price-to-Income między 2000 a 2024}
\begin{figure}[H]
    \raportimg{phts/1/png/1.7.png}
    \caption{Zmiana indeksu PTI krajów między 2000 a 2024 rokiem.}
\end{figure}

Wizualizacja pozwala określić, czy od roku 2000 do roku 2024 średnia względna wartość mieszkania wzrosła, czy zmalała.
Najbardziej odstającymi krajami są Luksemburg, gdzie wartość wzrosła, oraz Bułgaria, gdzie wartość zmalała.

\subsubsection{Zmiana Price-to-Income między minimum a maksimum}
\begin{figure}[H]
    \raportimg{phts/1/png/1.8.png}
    \caption{Zmiana indeksu PTI krajów między najmniejszą a największą dostępną wartością.}
\end{figure}

Na tym wykresie również dominuje Luksemburg, wyraźnie przewyższając pozostałe kraje.
Nie oznacza to jednak, że Luksemburg ma średnio najbardziej wartościowe mieszkania, lecz że doświadczył największej różnicy wartości indeksu w analizowanym okresie.

\subsubsection{Statystyki opisowe Rate of Change}
\begin{figure}[H]
    \raportimg{phts/1/png/1.6.png}
    \caption{Statystyki opisowe wartości indeksu RoC.}
\end{figure}

RoC wskazuje zmianę PTI względem poprzedniego roku.
Większość wartości RoC mieści się w przedziale od -10 pp. do +10 pp.
Średnia i mediana wskazują na niewielką przewagę dodatnich zmian.
W przedziale 2007--2012 widoczne są ujemne średnie wartości RoC, co można wiązać z kryzysem mieszkaniowym.

\subsubsection{Indeks Rate of Change względem kraju}
\begin{figure}[H]
    \raportimg{phts/1/png/1.3.png}
    \caption{Wartość indeksu RoC względem kraju na przestrzeni lat.}
    \label{wykres:1_8}
\end{figure}

RoC jest bardziej chaotyczną miarą niż sam PTI, ponieważ pokazuje zmiany rok do roku.
Rumunia ponownie wyróżnia się wartościami maksymalnymi, a w jednym z okresów wysokie wartości osiąga również Polska.

\subsubsection{Heatmap wartości Rate of Change}
\begin{figure}[H]
    \raportimg{phts/1/png/1.5.png}
    \caption{Heatmapa wartości Rate of Change krajów na przestrzeni lat.}
\end{figure}

Heatmapa pokazuje, jak kraje zmagały się ze zmianami indeksu PTI w czasie.
W latach 2005--2006 część krajów osiągała wysokie dodatnie wartości, które następnie spadały w latach 2008--2010.
Późniejszy okres charakteryzuje się bardziej mieszanymi zmianami, widocznymi również na wykresie Rate of Change.

\subsubsection{Wykresy dla Polski}
\begin{figure}[H]
    \raportimg{phts/1/png/1.10.png}
    \caption{Wykresy PTI i RoC na przestrzeni lat dla Polski.}
\end{figure}

W przypadku Polski widać spadek w latach 2008--2012, obecny również w wielu innych krajach Europy.
RoC pozostawało ujemne aż do 2017 roku, co oznacza, że wartości mieszkań względem dochodu statystycznie malały.
Od 2018 roku widoczny jest ponowny wzrost, utrzymujący się również w okresie pandemicznym, z wyjątkiem lat 2022 i 2023.

\subsection{Wnioski}
\begin{enumerate}
    \item W latach 2008--2015 w większości krajów Europy indeks PTI spadł, po czym stopniowo wzrastał.
    \item Większość wartości Rate of Change mieści się w relatywnie wąskim przedziale, co sugeruje umiarkowaną stabilność rynku mieszkaniowego.
    \item Specjalnym przypadkiem jest Rumunia, w której PTI i RoC są bardzo zmienne i odstające od pozostałych krajów.
\end{enumerate}

% ============================================================
\newpage
\section{Noclegi i miejsca do spania dla turystów w Europie}
% ============================================================

\subsection{Źródło danych}
\href{https://ec.europa.eu/eurostat/databrowser/view/tour_cap_nuts2dc__custom_12660565/default/table}
{Eurostat tour\_cap\_nuts2dc - Establishments and Bed-places in Tourist Accommodation by Urbanisation \& Coastal type}\\
\url{https://ec.europa.eu/eurostat/databrowser/view/tour_cap_nuts2dc__custom_12660565/default/table}

\subsection{Opis zbioru i istotnych pojęć}
\begin{itemize}
    \item Dane zostały zbierane na przestrzeni lat 2012--2024;
    \item Dane pochodzą z 27 krajów członkowskich Unii Europejskiej;
    \item Granulacja danych występuje względem państw, regionów NUTS2, stopnia urbanizacji oraz dostępu do morza;
    \item Dane dla niektórych krajów, zwłaszcza we wczesnych latach, są estymatami;
    \item \textbf{Placówka noclegowa / turystyczna} - ośrodek, który przyjmuje i kwateruje turystów z możliwością snu;
    \item \textbf{Miejsce do spania} - odpowiednie miejsce do odpoczynku nocnego dla jednej osoby.
\end{itemize}

\subsection{Wizualizacje zbioru}

\subsubsection{Kraje z największą liczbą placówek noclegowych}
\begin{figure}[H]
    \raportimg{phts/2/png/2.1.png}
    \caption{20 państw Europy z największą liczbą placówek noclegowych w 2024 roku.}
\end{figure}

Największą liczbą placówek noclegowych wyróżniają się Włochy, które znacząco przewyższają kolejne państwa.
Duża liczba placówek nie musi jednak oznaczać największej liczby dostępnych miejsc do spania, ponieważ placówki mogą mieć różną skalę działalności.

\subsubsection{Kraje z największą liczbą miejsc do spania}
\begin{figure}[H]
    \raportimg{phts/2/png/2.2.png}
    \caption{20 państw Europy z największą liczbą miejsc do spania w 2024 roku.}
\end{figure}

Włochy zajmują pierwsze miejsce również pod względem liczby miejsc do spania.
Chorwacja, mimo wysokiej liczby placówek, znajduje się niżej w rankingu miejsc do spania.
Pokazuje to, że nie każda placówka turystyczna ma taką samą liczbę łóżek.

\subsubsection{Średnia liczba miejsc do spania na placówkę turystyczną}
\begin{figure}[H]
    \raportimg{phts/2/png/2.7.png}
    \caption{Średnia liczba miejsc do spania na placówkę turystyczną.}
\end{figure}

Wykres pokazuje relację liczby miejsc do spania do liczby placówek.
Kraje z wysoką średnią mają relatywnie duże obiekty noclegowe, natomiast kraje z niską średnią mają więcej mniejszych placówek.
Włochy, mimo dominacji w liczbie placówek i miejsc, nie znajdują się na szczycie tej klasyfikacji.

\subsubsection{Zmienność liczby placówek noclegowych na przestrzeni lat}
\begin{figure}[H]
    \raportimg{phts/2/png/2.3.png}
    \caption{Wykres zmienności liczby placówek noclegowych na przestrzeni lat.}
\end{figure}

Dominacja Włoch i Chorwacji sprawia, że zmiany w pozostałych krajach są mniej widoczne.
Dlatego dodatkowo zastosowano wariant wykresu bez tych dwóch państw.

\begin{figure}[H]
    \raportimg{phts/2/png/2.11.png}
    \caption{Wykres zmienności liczby placówek noclegowych na przestrzeni lat -- bez Włoch i Chorwacji.}
\end{figure}

Po wyłączeniu Włoch i Chorwacji lepiej widać zmiany w pozostałych krajach, m.in. spadki i wzrosty liczby placówek w wybranych państwach.

\subsubsection{Indeks wzrostu liczby placówek turystycznych na przestrzeni lat}
\begin{figure}[H]
    \raportimg{phts/2/png/2.4.png}
    \caption{Wykres zmienności indeksu wzrostu liczby placówek turystycznych na przestrzeni lat.}
\end{figure}

Indeks wzrostu pozwala porównywać dynamikę zmian niezależnie od początkowej liczby placówek.
Dzięki temu kraje o mniejszej liczbie obiektów również mogą być porównywane z krajami większymi.

\subsubsection{Liczba placówek turystycznych w wybranych krajach, z podziałem na stopień urbanizacji}
\begin{figure}[H]
    \raportimg{phts/2/png/2.5.png}
    \caption{Liczba placówek turystycznych w wybranych krajach, z podziałem na stopień urbanizacji.}
\end{figure}

Wykres pokazuje, czy placówki koncentrują się przede wszystkim w miastach, na przedmieściach czy na obszarach wiejskich.
Pozwala to zauważyć różnice w charakterze turystyki w poszczególnych krajach.

\subsubsection{Liczba placówek turystycznych w wybranych krajach, z podziałem na dostęp do morza}
\begin{figure}[H]
    \raportimg{phts/2/png/2.6.png}
    \caption{Liczba placówek turystycznych w wybranych krajach, z podziałem na dostęp do morza.}
\end{figure}

Dostęp do morza jest ważnym czynnikiem dla wielu państw turystycznych.
Wykres pokazuje, w jakim stopniu baza noclegowa jest skoncentrowana na obszarach przybrzeżnych.

\subsubsection{Liczba placówek turystycznych i miejsc do spania w Polsce, z podziałem na województwa}
\begin{figure}[H]
    \raportimg{phts/2/png/2.8.png}
    \caption{Liczba placówek turystycznych i miejsc do spania w Polsce, z podziałem na województwa.}
\end{figure}

Wykres przedstawia zróżnicowanie bazy noclegowej w Polsce.
Największe wartości powinny pojawiać się w regionach silnie turystycznych oraz dużych ośrodkach miejskich.

\begin{figure}[H]
    \raportimg{phts/2/png/2.nuts2.png}
    \caption{NUTS2 -- podział Polski na wybrane jednostki geograficzne.}
\end{figure}

\subsubsection{Liczba placówek turystycznych i miejsc do spania w Polsce, z podziałem na stopień urbanizacji oraz dostęp do morza}
\begin{figure}[H]
    \raportimg{phts/2/png/2.9.png}
    \caption{Liczba placówek turystycznych i miejsc do spania w Polsce, z podziałem na stopień urbanizacji oraz dostęp do morza.}
\end{figure}

Wizualizacja pokazuje, jak w Polsce baza noclegowa rozkłada się względem urbanizacji oraz położenia przybrzeżnego.

\subsubsection{Dodatkowe wizualizacje dla Polski}
\begin{figure}[H]
    \raportimg{phts/2/png/2.10.png}
    \caption{Dodatkowe wizualizacje dla Polski.}
\end{figure}

\subsection{Wnioski}
\begin{enumerate}
    \item Włochy dominują pod względem liczby placówek noclegowych oraz liczby miejsc do spania.
    \item Liczba placówek nie przekłada się bezpośrednio na liczbę miejsc do spania, ponieważ placówki różnią się wielkością.
    \item Analiza z wyłączeniem krajów dominujących pozwala lepiej zauważyć zmiany w pozostałych państwach.
    \item Dostęp do morza oraz stopień urbanizacji są istotnymi czynnikami różnicującymi rozmieszczenie bazy noclegowej.
\end{enumerate}

% ============================================================
\newpage
\section{Nie-dominujące wyznania religijne w Polsce}
% ============================================================

\subsection{Źródło danych}
% \href{}{}\\
% \url{}

\subsection{Opis zbioru i istotnych pojęć}
\begin{itemize}
    \item
\end{itemize}

\subsection{Wizualizacje zbioru}

\subsection{Wnioski}

% ============================================================
\newpage
\section{Kwartalny dług publiczny}
% ============================================================

\subsection{Źródło danych}
\href{https://ec.europa.eu/eurostat/databrowser/view/gov_10q_ggdebt/default/table?lang=en}
{Eurostat gov\_10q\_ggdebt - Quarterly government debt}\\
\url{https://ec.europa.eu/eurostat/databrowser/view/gov_10q_ggdebt/default/table?lang=en}

\subsection{Opis zbioru i istotnych pojęć}
\begin{itemize}
    \item Zbiór przedstawia kwartalne dane dotyczące długu sektora instytucji rządowych i samorządowych;
    \item Dane zostały zebrane na przestrzeni lat 2012--2025;
    \item Dane pochodzą z krajów członkowskich Unii Europejskiej;
    \item Zbiór analizowany w prezentacji pochodzi z Eurostatu i ma oznaczenie \texttt{gov\_10q\_ggdebt};
    \item Ostatnia aktualizacja danych w wykorzystanym zbiorze miała miejsce 22/01/2026;
    \item \textbf{Dług publiczny} - zadłużenie sektora instytucji rządowych i samorządowych;
    \item \textbf{Indeks długu} - wartość długu przeliczona względem pierwszej dostępnej obserwacji;
    \item \textbf{Dynamika r/r} - procentowa zmiana wartości długu względem analogicznego kwartału poprzedniego roku;
    \item \textbf{Odchylenie standardowe dynamiki r/r} - miara niestabilności zmian zadłużenia w czasie.
\end{itemize}

\subsection{Wizualizacje zbioru}

\subsubsection{Indeks długu publicznego dla 8 krajów o największym wzroście}
\begin{figure}[H]
    \raportimg{phts/4/png/4.1.png}
    \caption{Dług publiczny: indeks 2012-Q1 = 100 dla 8 krajów o największym wzroście.}
\end{figure}

Wykres przedstawia kraje, w których indeks długu publicznego wzrósł najmocniej względem początku analizowanego okresu.
Najbardziej wyróżnia się Estonia, a silny wzrost widoczny jest również w Rumunii oraz Bułgarii, szczególnie po roku 2020.

\subsubsection{Skumulowana zmiana długu od początku próby}
\begin{figure}[H]
    \raportimg{phts/4/png/4.2.png}
    \caption{Skumulowana zmiana długu publicznego od początku próby.}
\end{figure}

Najsilniejszy wzrost odnotowano w Estonii, gdzie skumulowana zmiana wyniosła około 586\%.
Wysokie wartości występują również w Rumunii oraz Bułgarii.
Polska znajduje się w grupie krajów z wyraźnym wzrostem długu, jednak nie osiąga wartości tak wysokich jak kraje dominujące.

\subsubsection{Indeks długu publicznego dla wybranych dużych gospodarek}
\begin{figure}[H]
    \raportimg{phts/4/png/4.3.png}
    \caption{Dług publiczny: indeks 2012-Q1 = 100 dla wybranych dużych gospodarek.}
\end{figure}

Wykres porównuje indeks długu publicznego dla Niemiec, Hiszpanii, Francji, Włoch oraz Polski.
Największy wzrost w tej grupie widoczny jest dla Polski, szczególnie po roku 2020.

\subsubsection{Dynamika długu rok do roku dla krajów o najsilniejszych bieżących zmianach}
\begin{figure}[H]
    \raportimg{phts/4/png/4.4.png}
    \caption{Dynamika długu rok do roku -- 5 krajów o najsilniejszych bieżących zmianach.}
\end{figure}

Dynamika długu jest bardziej chaotyczna niż sam indeks długu.
Widać krótkookresowe przyspieszenia i spowolnienia zmian zadłużenia.

\subsubsection{Najwyższa zmienność dynamiki długu rok do roku}
\begin{figure}[H]
    \raportimg{phts/4/png/4.5.png}
    \caption{Najwyższa zmienność dynamiki długu publicznego rok do roku -- top 10 krajów.}
\end{figure}

Największą zmienność wykazuje Estonia, co oznacza, że tempo zmian jej długu publicznego było najmniej stabilne spośród analizowanych państw.

\subsubsection{Heatmap indeksu długu publicznego}
\begin{figure}[H]
    \raportimg{phts/4/png/4.6.png}
    \caption{Heatmapa indeksu długu publicznego względem 2012-Q1.}
\end{figure}

Heatmapa przedstawia przebieg indeksu długu publicznego dla wszystkich analizowanych krajów.
Kolory pozwalają szybko zauważyć, które kraje utrzymywały się blisko punktu startowego, a które weszły na wyraźnie wyższe poziomy zadłużenia.

\subsubsection{Ostatni kwartał względem średniej z całego okresu}
\begin{figure}[H]
    \raportimg{phts/4/png/4.7.png}
    \caption{Ostatni kwartał względem średniej z całego okresu -- top 10 krajów.}
\end{figure}

Największe dodatnie odchylenie występuje w Rumunii, Estonii oraz Bułgarii.
Oznacza to, że obecny poziom zadłużenia tych państw jest wyraźnie wyższy niż ich przeciętny poziom z lat 2012--2025.

\subsubsection{Średni indeks długu w grupach krajów}
\begin{figure}[H]
    \raportimg{phts/4/png/4.8.png}
    \caption{Średni indeks długu publicznego w grupach krajów.}
\end{figure}

Największy wzrost średniego indeksu widoczny jest w grupie państw Europy Środkowo-Wschodniej.
Po 2020 roku krzywa dla tej grupy zaczyna rosnąć znacznie szybciej.

\subsection{Wnioski}
\begin{enumerate}
    \item Najsilniejszy wzrost nominalnego długu od początku próby odnotowano w Estonii.
    \item Największą zmienność dynamiki rok do roku miała Estonia.
    \item Po 2020 roku w wielu krajach widoczne jest przyspieszenie wzrostu długu publicznego.
    \item Polska, w porównaniu z wybranymi dużymi gospodarkami europejskimi, charakteryzuje się jednym z mocniejszych wzrostów indeksu długu względem 2012-Q1.
\end{enumerate}

% ============================================================
\newpage
\section{Misje kosmiczne}
% ============================================================

\subsection{Źródło danych}
\href{https://www.kaggle.com/datasets/agirlcoding/all-space-missions-from-1957}
{Kaggle - All Space Missions From 1957}\\
\url{https://www.kaggle.com/datasets/agirlcoding/all-space-missions-from-1957}

\subsection{Opis zbioru i istotnych pojęć}
\begin{itemize}
    \item Zbiór przedstawia globalne misje kosmiczne realizowane przez różne kraje, firmy oraz organizacje;
    \item Dane obejmują okres od 1957 do 2020 roku;
    \item Ostatnia aktualizacja wykorzystanego zbioru miała miejsce 13/08/2020;
    \item Dane kosztowe są dostępne tylko dla części misji;
    \item Dekada 2020 nie jest pełna, dlatego została pominięta w analizie dekad;
    \item \textbf{Misja kosmiczna} - pojedyncze przedsięwzięcie związane ze startem rakiety lub pojazdu kosmicznego;
    \item \textbf{Status misji} - końcowy wynik misji;
    \item \textbf{Skuteczność misji} - odsetek misji zakończonych sukcesem względem wszystkich misji w danej grupie.
\end{itemize}

\subsection{Wizualizacje zbioru}

\subsubsection{Liczba misji kosmicznych w czasie}
\begin{figure}[H]
    \raportimg{phts/5/png/5.1.png}
    \caption{Liczba misji kosmicznych w czasie.}
\end{figure}

Największa aktywność widoczna jest w okresie zimnej wojny, szczególnie od drugiej połowy lat 60. do końca lat 70.
Ponowny wzrost aktywności widoczny jest w drugiej połowie lat 2010.

\subsubsection{Liczba misji kosmicznych według dekad}
\begin{figure}[H]
    \raportimg{phts/5/png/5.2.png}
    \caption{Liczba misji kosmicznych według dekad.}
\end{figure}

Najwięcej misji przeprowadzono w latach 70.
Po tym okresie nastąpił spadek liczby misji, a w latach 2010. ponowny wzrost aktywności sektora kosmicznego.

\subsubsection{Top 10 krajów według liczby misji}
\begin{figure}[H]
    \raportimg{phts/5/png/5.3.png}
    \caption{Top 10 krajów według liczby misji kosmicznych według lokalizacji startu.}
\end{figure}

Najwięcej misji przeprowadzono z Rosji oraz Stanów Zjednoczonych.
Na trzecim miejscu znajduje się Kazachstan, głównie ze względu na kosmodrom Bajkonur.

\subsubsection{Top 10 firm i organizacji według liczby misji}
\begin{figure}[H]
    \raportimg{phts/5/png/5.4.png}
    \caption{Top 10 firm i organizacji według liczby misji kosmicznych.}
\end{figure}

Dominuje RVSN USSR, co pokazuje skalę radzieckiej aktywności kosmicznej w analizowanym okresie.

\subsubsection{Aktywność top 5 firm w czasie}
\begin{figure}[H]
    \raportimg{phts/5/png/5.5.png}
    \caption{Aktywność top 5 firm i organizacji w czasie.}
\end{figure}

Najbardziej charakterystyczny jest przebieg RVSN USSR, które osiągało bardzo wysoką liczbę misji w latach 60. i 70.
Szczególnie widoczny jest również wzrost aktywności CASC w końcowej części szeregu czasowego, co wskazuje na rosnące znaczenie Chin w sektorze kosmicznym.

\subsubsection{Sezonowość startów kosmicznych według miesięcy}
\begin{figure}[H]
    \raportimg{phts/5/png/5.6.png}
    \caption{Sezonowość startów kosmicznych według miesięcy.}
\end{figure}

Najwięcej startów odnotowano w grudniu, a wysokie wartości widoczne są również w czerwcu i październiku.
Najmniejsza liczba misji przypada na styczeń.

\subsubsection{Rozkład wyników misji}
\begin{figure}[H]
    \raportimg{phts/5/png/5.7.png}
    \caption{Rozkład wyników misji kosmicznych.}
\end{figure}

Zdecydowana większość misji zakończyła się sukcesem.
Porażki, częściowe porażki oraz porażki przedstartowe stanowią znacznie mniejszą część zbioru.

\subsubsection{Skuteczność misji w czasie}
\begin{figure}[H]
    \raportimg{phts/5/png/5.8.png}
    \caption{Skuteczność misji kosmicznych w czasie.}
\end{figure}

Na początku ery kosmicznej skuteczność była bardziej niestabilna.
W późniejszych latach ustabilizowała się na wysokim poziomie, najczęściej w okolicach 90\% lub więcej.

\subsubsection{Skuteczność misji kosmicznych według dekad}
\begin{figure}[H]
    \raportimg{phts/5/png/5.9.png}
    \caption{Skuteczność misji kosmicznych według dekad.}
\end{figure}

Skuteczność misji wzrosła znacząco od lat 50. do kolejnych dekad.
Od lat 70. utrzymuje się na bardzo wysokim poziomie.

\subsubsection{Skuteczność misji w top 10 krajów}
\begin{figure}[H]
    \raportimg{phts/5/png/5.10.png}
    \caption{Skuteczność misji w top 10 krajów według lokalizacji startu.}
\end{figure}

Wśród krajów z dużą liczbą misji wysoką skuteczność osiągają m.in. Francja, Rosja, Japonia oraz Chiny.
Wyniki krajów z bardzo małą liczbą misji należy interpretować ostrożnie.

\subsubsection{Koszt misji kosmicznych w czasie}
\begin{figure}[H]
    \raportimg{phts/5/png/5.11.png}
    \caption{Koszt misji kosmicznych w czasie.}
\end{figure}

Widoczna jest ogólna tendencja spadkowa kosztów misji.
Należy jednak pamiętać, że dane kosztowe są dostępne tylko dla części obserwacji.

\subsection{Wnioski}
\begin{enumerate}
    \item Największa liczba misji kosmicznych przypadła na lata 60. i 70.
    \item Najwięcej startów przeprowadzono z Rosji oraz Stanów Zjednoczonych.
    \item Najaktywniejszą organizacją w całym zbiorze było RVSN USSR.
    \item Skuteczność misji kosmicznych wzrosła wraz z rozwojem technologii.
    \item W ostatnich latach widoczny jest ponowny wzrost liczby misji.
    \item Koszt misji kosmicznych wykazuje tendencję spadkową, ale analiza kosztów dotyczy tylko części obserwacji.
\end{enumerate}

% ============================================================
\newpage
\section{Temperatura i opady w wybranych miastach Europy (R)}

\subsection{Źródło danych}
Dane wykorzystane w tej części projektu pochodzą z czterech plików CSV przygotowanych dla wybranych miast europejskich: Warszawy, Paryża, Dublina oraz Lizbony.

\begin{itemize}
    \item \texttt{warsaw.csv} -- dane dla Warszawy;
    \item \texttt{paris.csv} -- dane dla Paryża;
    \item \texttt{dublin.csv} -- dane dla Dublina;
    \item \texttt{lisbona.csv} -- dane dla Lizbony.
\end{itemize}

\subsection{Opis zbioru i istotnych pojęć}
\begin{itemize}
    \item Zbiór przedstawia miesięczne dane klimatyczne dla czterech miast: Warszawy, Paryża, Dublina i Lizbony;
    \item Dane obejmują różne zakresy czasowe w zależności od miasta: Warszawa od 2007-12 do 2026-03, Paryż od 2003-03 do 2026-03, Dublin od 2000-01 do 2026-03 oraz Lizbona od 2000-01 do 2026-03;
    \item Każda obserwacja odpowiada jednemu miesiącowi w danym mieście;
    \item \textbf{Temperatura} oznacza średnią miesięczną temperaturę powietrza wyrażoną w stopniach Celsjusza;
    \item \textbf{Opady} oznaczają miesięczny poziom opadów wyrażony w milimetrach;
    \item \textbf{Pora roku} została wyznaczona na podstawie miesiąca: wiosna obejmuje marzec, kwiecień i maj; lato czerwiec, lipiec i sierpień; jesień wrzesień, październik i listopad; zima grudzień, styczeń i luty;
    \item \textbf{Średnia krocząca 12-miesięczna} pozwala wygładzić miesięczne wahania temperatury i lepiej pokazać długookresowy trend;
    \item \textbf{Anomalia temperatury} oznacza różnicę między rzeczywistą temperaturą w danym miesiącu a średnią temperaturą typową dla tego samego miesiąca w całym analizowanym okresie;
    \item \textbf{Trend odporny} na wykresie zależności temperatury i opadów został zastosowany po to, aby pojedyncze skrajne obserwacje opadowe nie zniekształcały nadmiernie linii trendu.
\end{itemize}

\subsection{Wizualizacje zbioru}

\subsubsection{Warszawa}

\wykrespngwniosek{phts/6/png/6.1.png}{Warszawa: temperatura miesięczna i trend.}{wykres-6-1}{Warszawa ma bardzo wyraźną sezonowość temperatur: zimą wartości spadają często w okolice zera lub poniżej, a latem regularnie przekraczają 20\textdegree C. Linia średniej kroczącej pokazuje stopniowy wzrost temperatur w późniejszej części szeregu czasowego.}

\wykrespngwniosek{phts/6/png/6.2.png}{Warszawa: średnia roczna temperatura.}{wykres-6-2}{W kolejnych latach coraz częściej pojawiają się wartości powyżej średniej długookresowej. Widoczny jest ogólny wzrost średniej rocznej temperatury, choć pojedyncze niepełne lata na początku i końcu szeregu należy interpretować ostrożnie.}

\wykrespngwniosek{phts/6/png/6.3.png}{Warszawa: anomalie temperatury.}{wykres-6-3}{W pierwszej części szeregu częściej występują miesiące chłodniejsze od normy, natomiast w późniejszym okresie częściej pojawiają się dodatnie anomalie. Sugeruje to przesunięcie w stronę cieplejszych miesięcy względem typowych wartości sezonowych.}

\wykrespngwniosek{phts/6/png/6.4.png}{Warszawa: typowy przebieg temperatur w roku.}{wykres-6-4}{Roczny cykl temperatur jest bardzo silny: najniższe wartości przypadają na styczeń i luty, a najwyższe na lipiec i sierpień. Taki przebieg wskazuje na bardziej kontynentalny charakter klimatu Warszawy.}

\wykrespngwniosek{phts/6/png/6.5.png}{Warszawa: rozkład temperatur w porach roku.}{wykres-6-5}{Największy kontrast występuje między zimą i latem. Rozkład zimowy obejmuje wartości ujemne, natomiast lato jest wyraźnie przesunięte ku wysokim temperaturom, co potwierdza dużą amplitudę sezonową.}

\wykrespngwniosek{phts/6/png/6.6.png}{Warszawa: średnie opady miesięczne.}{wykres-6-6}{Najwyższe średnie opady przypadają na okres letni, szczególnie na lipiec. Oznacza to, że w Warszawie cieplejsza część roku jest jednocześnie okresem większych opadów.}

\wykrespngwniosek{phts/6/png/6.7.png}{Warszawa: suma opadów w latach.}{wykres-6-7}{Roczne sumy opadów są mocno zróżnicowane. Pojawiają się zarówno lata wyraźnie wilgotniejsze od średniej, jak i lata suche, co wskazuje na dużą zmienność opadów między latami.}

\wykrespngwniosek{phts/6/png/6.8.png}{Warszawa: temperatura a opady.}{wykres-6-8}{Widoczna jest dodatnia zależność między temperaturą a opadami, ale wynika ona głównie z sezonowości. Latem temperatury są wysokie i jednocześnie częściej występują większe opady. Pojedyncze bardzo wysokie opady nie powinny być usuwane, lecz wymagają ostrożnej interpretacji.}

\FloatBarrier
\subsubsection{Paryż}

\wykrespngwniosek{phts/6/png/6.9.png}{Paryż: temperatura miesięczna i trend.}{wykres-6-9}{Paryż ma wyraźną sezonowość temperatur, ale łagodniejszą niż Warszawa. Trend jest mniej jednoznaczny, a wahania miesięczne są umiarkowane.}

\wykrespngwniosek{phts/6/png/6.10.png}{Paryż: średnia roczna temperatura.}{wykres-6-10}{Średnie roczne temperatury oscylują wokół średniej długookresowej. W ostatnich latach częściej pojawiają się wartości powyżej średniej, jednak zmiany nie są tak wyraźne jak w Warszawie.}

\wykrespngwniosek{phts/6/png/6.11.png}{Paryż: anomalie temperatury.}{wykres-6-11}{Anomalie pokazują okresy zarówno chłodniejsze, jak i cieplejsze od normy miesięcznej. W końcowej części szeregu widoczna jest większa liczba dodatnich anomalii, co może wskazywać na częstsze miesiące cieplejsze niż typowo.}

\wykrespngwniosek{phts/6/png/6.12.png}{Paryż: typowy przebieg temperatur w roku.}{wykres-6-12}{Najcieplejsze miesiące to lipiec i sierpień, a najchłodniejsze styczeń i luty. Roczny przebieg temperatur jest regularny, ale różnice między zimą i latem są mniejsze niż w Warszawie.}

\wykrespngwniosek{phts/6/png/6.13.png}{Paryż: rozkład temperatur w porach roku.}{wykres-6-13}{Rozkłady sezonowe są dobrze rozdzielone, ale zima jest łagodniejsza niż w Warszawie. Lato osiąga wysokie temperatury, natomiast wiosna i jesień stanowią wyraźne okresy przejściowe.}

\wykrespngwniosek{phts/6/png/6.14.png}{Paryż: średnie opady miesięczne.}{wykres-6-14}{Opady w Paryżu są dość równomiernie rozłożone w ciągu roku. Nie występuje tak silnie zaznaczona pora sucha jak w Lizbonie ani tak wyraźne letnie maksimum jak w Warszawie.}

\wykrespngwniosek{phts/6/png/6.15.png}{Paryż: suma opadów w latach.}{wykres-6-15}{Roczne sumy opadów zmieniają się między latami, ale większość wartości pozostaje w umiarkowanym zakresie. Pojedyncze lata wilgotniejsze od średniej wyróżniają się, lecz nie dominują całego szeregu.}

\wykrespngwniosek{phts/6/png/6.16.png}{Paryż: temperatura a opady.}{wykres-6-16}{Zależność między temperaturą i opadami jest słaba. Punkty są rozproszone, a linia trendu ma niewielkie nachylenie, co oznacza, że w Paryżu sama temperatura nie wyjaśnia dobrze poziomu miesięcznych opadów.}

\FloatBarrier
\subsubsection{Dublin}

\wykrespngwniosek{phts/6/png/6.17.png}{Dublin: temperatura miesięczna i trend.}{wykres-6-17}{Dublin ma najmniejszą sezonowość temperatur spośród analizowanych miast. Temperatury są stosunkowo łagodne przez cały rok, a linia trendu pozostaje dość stabilna.}

\wykrespngwniosek{phts/6/png/6.18.png}{Dublin: średnia roczna temperatura.}{wykres-6-18}{Średnie roczne temperatury są skupione blisko średniej długookresowej. Oznacza to relatywnie stabilny przebieg temperatur w czasie, choć ostatnie pełne lata częściej przekraczają średnią.}

\wykrespngwniosek{phts/6/png/6.19.png}{Dublin: anomalie temperatury.}{wykres-6-19}{Anomalie temperatury są przeważnie mniejsze niż w Warszawie, co potwierdza bardziej stabilny charakter klimatu. W końcowej części szeregu widocznych jest więcej dodatnich odchyleń od normy miesięcznej.}

\wykrespngwniosek{phts/6/png/6.20.png}{Dublin: typowy przebieg temperatur w roku.}{wykres-6-20}{Różnica między zimą i latem jest stosunkowo niewielka. Najcieplejszy okres przypada na lipiec i sierpień, ale wartości letnie są niższe niż w pozostałych analizowanych miastach.}

\wykrespngwniosek{phts/6/png/6.21.png}{Dublin: rozkład temperatur w porach roku.}{wykres-6-21}{Rozkłady temperatur w porach roku są do siebie bardziej zbliżone niż w Warszawie czy Paryżu. Zima jest łagodna, a lato umiarkowanie ciepłe, co wskazuje na silny wpływ klimatu oceanicznego.}

\wykrespngwniosek{phts/6/png/6.22.png}{Dublin: średnie opady miesięczne.}{wykres-6-22}{Opady utrzymują się na wysokim poziomie przez większość roku. Największe wartości przypadają na jesień, zwłaszcza październik i listopad.}

\wykrespngwniosek{phts/6/png/6.23.png}{Dublin: suma opadów w latach.}{wykres-6-23}{Roczne sumy opadów są dość wysokie, ale niektóre lata wyraźnie odbiegają od średniej. Ogólnie Dublin pozostaje miastem wilgotnym w porównaniu z Warszawą i Paryżem.}

\wykrespngwniosek{phts/6/png/6.24.png}{Dublin: temperatura a opady.}{wykres-6-24}{Związek między temperaturą i opadami jest bardzo słaby. Opady występują w różnych zakresach temperatur, co oznacza, że w Dublinie wilgotność nie jest ograniczona tylko do jednej pory roku.}

\FloatBarrier
\subsubsection{Lizbona}

\wykrespngwniosek{phts/6/png/6.25.png}{Lizbona: temperatura miesięczna i trend.}{wykres-6-25}{Lizbona jest najcieplejszym z analizowanych miast. Temperatury pozostają wysokie przez cały rok, a zimy są łagodne. Trend wskazuje na utrzymywanie się wysokich temperatur w ostatnich latach.}

\wykrespngwniosek{phts/6/png/6.26.png}{Lizbona: średnia roczna temperatura.}{wykres-6-26}{Średnie roczne temperatury są najwyższe spośród czterech miast. W wielu latach wartości przekraczają średnią długookresową, choć niepełne lata na początku i końcu szeregu należy traktować ostrożnie.}

\wykrespngwniosek{phts/6/png/6.27.png}{Lizbona: anomalie temperatury.}{wykres-6-27}{Anomalie wskazują, że w późniejszej części szeregu częściej pojawiają się miesiące cieplejsze od normy. Jednocześnie występują pojedyncze silne ujemne anomalie, które pokazują, że nawet w ciepłym klimacie możliwe są chłodniejsze epizody.}

\wykrespngwniosek{phts/6/png/6.28.png}{Lizbona: typowy przebieg temperatur w roku.}{wykres-6-28}{Roczny przebieg temperatur pokazuje bardzo łagodne zimy oraz gorące lato. Najwyższe wartości przypadają na lipiec i sierpień, a nawet zimą średnie temperatury pozostają stosunkowo wysokie.}

\wykrespngwniosek{phts/6/png/6.29.png}{Lizbona: rozkład temperatur w porach roku.}{wykres-6-29}{Wszystkie pory roku są cieplejsze niż w pozostałych miastach. Szczególnie wysoko położony jest rozkład letni, natomiast zima nadal pozostaje łagodna.}

\wykrespngwniosek{phts/6/png/6.30.png}{Lizbona: średnie opady miesięczne.}{wykres-6-30}{Rozkład opadów jest bardzo nierównomierny. Lato jest prawie suche, zwłaszcza lipiec i sierpień, natomiast największe opady występują jesienią i zimą.}

\wykrespngwniosek{phts/6/png/6.31.png}{Lizbona: suma opadów w latach.}{wykres-6-31}{Roczne sumy opadów są bardzo zmienne. Występują zarówno lata wyjątkowo suche, jak i bardzo wilgotne, co jest typowe dla klimatu śródziemnomorskiego.}

\wykrespngwniosek{phts/6/png/6.32.png}{Lizbona: temperatura a opady.}{wykres-6-32}{Widoczna jest wyraźna zależność ujemna: cieplejsze miesiące są zwykle bardziej suche, a chłodniejsze miesiące częściej wiążą się z większymi opadami. Jest to typowy wzorzec klimatu śródziemnomorskiego.}

\FloatBarrier
\subsection{Wnioski ogólne}
\begin{enumerate}
    \item Spośród analizowanych miast najwyższe temperatury występują w Lizbonie, natomiast największa sezonowość temperatur widoczna jest w Warszawie.
    \item Dublin ma najbardziej wyrównany przebieg temperatur w ciągu roku, co potwierdza wpływ klimatu oceanicznego.
    \item Paryż zajmuje pozycję pośrednią między klimatem bardziej kontynentalnym a oceanicznym.
    \item Lizbona charakteryzuje się wyraźnie suchym latem i większymi opadami w chłodniejszej części roku.
    \item Analiza anomalii temperatury pozwala lepiej wychwycić miesiące nietypowo ciepłe i chłodne niż bezpośrednia analiza samych temperatur.
    \item Zależność temperatury i opadów jest częściowo kształtowana przez sezonowość i różni się między miastami.
    \item Ujęcie oddzielnych wykresów dla każdego miasta poprawia czytelność analizy i pozwala zachować pełną numerację rysunków od 6.1 do 6.32.
\end{enumerate}

\end{document}
